\section{Discussion}
\label{sec:discussion}

\paragraph{Why a language model beats the models we usually tune.}
A GBDT on 21 columns learns a function of the marginals and a handful of interactions it
happens to split on. The data, however, were drawn from a joint distribution over 21
columns with strong dependence: delays, distance and cabin co-vary, and the generator
reproduces that joint at a temperature low enough to sharpen it. Scoring $p(x\mid y)$ uses
all of it, including the parts that do not predict the label but do constrain which rows
are plausible. The gap between 0.961 (discriminative GPT-2) and 0.9624 (same network, other
direction) is the value of that joint structure, not of scale: GPT-2 medium and large are
no better than small, while the direction is worth more than twenty times a new GBDT.

\paragraph{Is it a leak?}
No, and the write-ups are careful about it. The original survey is a public dataset that
the competition rules allow; fitting a density to it and applying Bayes' rule is a modelling
choice. The 4th-place result --- 0.96201 from a transformer trained from scratch on the
competition rows, with learned embeddings and no GPT-2 weights --- shows that the direction
matters even without the original corpus, and the 1st-place team re-ran the target-last
recipe on non-synthetic data and still beat XGBoost. What the generator gives is a prior
over rows; the competition still decides who estimates it best.

\paragraph{Why the stacks converge to 0.9624.}
Public OOF predictions are not independent evidence: they are many estimators of the same
conditional, trained on the same 700k rows, with the same features. A logistic head over
250 of them behaves like a very well-regularised average, and its ceiling is the ceiling of
that conditional. This is also why the private board ties at five decimals for 126 teams:
the differences that remain are the arbitrary ones.

\paragraph{Limitations of this survey.}
(i) We did not rerun any winning pipeline; every number in Sections~\ref{sec:generator}
and~\ref{sec:evidence} is self-reported, and the baselines differ (5 folds with different
seeds, different target encodings). (ii) Five of the top ten published nothing; we name
their reported backbones only where other write-ups credit them. (iii) Stack scores on
pooled public OOF are mildly optimistic, as the 25th-place team states explicitly: the
members were fitted on the same labels the stacker is scored against. (iv) The private
scores we use come from the leaderboard, whose ties at five decimals hide differences below
$10^{-5}$.

\paragraph{What generalises beyond this competition.}
Two practices transfer to any tabular problem. First, \emph{identify the data-generating
process before tuning}: here a tokenizer fingerprint and a temperature sweep cost an
afternoon and were worth \dlt{140} over the best conventional pipeline. Second, \emph{make
the validation honest enough to prefer a null result}: a frozen split, a file contract,
nested stacking and a pre-registered screen are cheap, and they are what separated the
teams that climbed from the teams that fell.

For our own repository, the actionable list is short: treat \texttt{flight\_distance} as a
route key with fold-safe target statistics (worth \dlt{15} on a single LightGBM in the
29th-place measurement), add a generative scorer as a member rather than as a replacement,
and re-tune the stacker's shrinkage whenever the pool grows.

\section{Reproducing this survey}
\label{sec:repro}

\begin{sloppypar}
Everything used here is public and was collected with the Kaggle CLI into
\texttt{artifacts/solutions/}: \texttt{leaderboard\_full.csv} (public board, 3{,}735 rows),
\texttt{leaderboard\_top200.csv} (private top 200), the ten write-ups as plain text under
\texttt{discussions/}, and one public stacking notebook. The two plotted figures are
produced by \texttt{paper/solutions/figures.py}; the two diagrams are TikZ in this
document.
\end{sloppypar}

\begin{verbatim}
kaggle competitions leaderboard playground-series-s6e10 -d
kaggle competitions leaderboard playground-series-s6e10 --show -v --page-size 200
kaggle competitions topics list playground-series-s6e10 --sort-by top
kaggle competitions topics show playground-series-s6e10 754683
python paper/solutions/figures.py
\end{verbatim}

Two practical notes for Windows users: the leaderboard archive contains a colon in its
entry name, so \texttt{Expand-Archive} refuses it and the CSV has to be extracted with
Python's \texttt{zipfile}; and the CLI prints Unicode minus signs, which requires
\texttt{PYTHONIOENCODING=utf-8} in a GBK console.

\appendix
\setlength{\emergencystretch}{2em}
\section{Recipe card for a generative scorer}
\label{app:recipe}

The following is the minimal version of the winning recipe, assembled from the four
published implementations. It assumes only a causal language model, no feature tokenizer
and no GPU cluster: GPT-2 small fits in 8\,GB and the five folds take a few hours on a
single mid-range GPU.

\begin{verbatim}
# 1. serialise a row exactly like the original file
text(row, y) = "satisfaction is {y}, " + ", ".join(f"{c} is {v}" for c, v in row)
               # integer age, one-decimal delays, "None" for missing

# 2. fit the generator on the original survey (129,880 rows)
model = GPT2small()
train_lm(model, [text(row, y) for row, y in original], epochs=60)

# 3. per fold: fine-tune, then score every validation and test row twice
for k in range(5):
    fold_model = train_lm(clone(model),
                          [text(row, y) for row, y in train if fold_id[row] != k],
                          epochs=8)
    for row in valid[k] + test:
        for y in (True, False):
            L[y] = mean_over_orders(sum_j log p_fold_model(x_j | x_<j, y))
        z[row] = log(pi / (1 - pi)) + L[True] - L[False]
\end{verbatim}

Four details decide whether this reproduces 0.9624 or 0.9610: (i) the label must come
\emph{first}; (ii) the serialisation must match the file the generator was trained on
(integer vs.\ float formatting, \texttt{None} for missing values); (iii) use a tokenizer
that keeps numbers compact --- GPT-2's, not a digit-wise one; and (iv) average the
log-likelihood difference over several column orders (8--32), which removes the arbitrary
order the model would otherwise memorise. The 7th-place team's variant replaces steps 2--3
with a from-scratch 6-layer transformer over per-column token vocabularies and keeps the
same scoring rule; the 4th-place variant does the same with 32 orders and a learned column
embedding. Both reach the same region, which is the strongest evidence that the direction,
not the backbone, is doing the work.

\sloppy
\bibliographystyle{plainnat}
\bibliography{references}
